\chapter{Research Objectives}\label{sec:research-objectives}

\section{Main Objective}\label{sec:main-objective}

To systematically evaluate and compare classical and AI-based relay strategies for two-hop communication, and to determine the optimal relay architecture as a function of channel conditions, SNR regime, and computational constraints.

\section{Research Gap and Motivation}\label{sec:research-gap-and-motivation}

Despite growing interest in AI for wireless communication, a systematic comparison of relay processing paradigms : spanning classical signal processing, supervised learning, generative modeling, and modern sequence architectures : is absent from the literature. 

This thesis addresses all five gaps through a unified framework that implements and compares nine relay strategies across six channel/topology configurations, with both original and normalized parameter counts, full statistical analysis, and a dedicated complexity study.



The following specific gaps motivate this thesis

\begin{enumerate}
    \item \textbf{Gap 1: No cross-paradigm relay comparison.} Existing work on AI-based relay processing focuses on individual architectures in isolation. Ye et al.~\cite{YeLiJuang2018OFDMDL} demonstrated DNNs for OFDM detection, Dorner et al.~\cite{Dorner2018OverTheAirDL} explored autoencoders for end-to-end communication, and Samuel et al.~\cite{SamuelDiskinWunder2019MIMODetect} applied unfolded networks to MIMO detection. However, no study systematically compares supervised feedforward networks, variational autoencoders, adversarial networks, attention-based Transformers, and state space models (Mamba) for the same relay denoising task under controlled conditions. Without such a comparison, practitioners cannot make informed architecture selections.
    \item \textbf{Gap 2: Model complexity-performance relationship is unknown.} It is commonly assumed that larger neural networks yield better performance. However, the relay denoising task : recovering a BPSK symbol from a noisy observation : is a low-dimensional mapping problem. The theoretical minimum description length for this mapping is small (a soft threshold function), suggesting that large models may overfit. No prior work has rigorously characterized the complexity-performance trade-off for relay processing, nor has a normalized (equal-parameter) comparison been conducted to isolate architectural effects from capacity effects.

    \item \textbf{Gap 3: Limited channel diversity in evaluations.} Most AI-for-wireless studies evaluate on a single channel model (typically AWGN or Rayleigh). The performance of AI relays under Rician fading, MIMO spatial multiplexing, and different equalization strategies has not been studied. Whether AI relays trained on one channel generalize to others : and whether their advantage over classical methods persists across diverse propagation conditions : are open questions.

    \item \textbf{Gap 4: State space models for physical-layer signal processing.} The Mamba architecture \cite{GuGoelRe2022SSM} and its Mamba-2 SSD variant \cite{DaoGu2024TransformersSSM} have demonstrated strong performance in NLP and genomics, but their application to physical-layer wireless communication is entirely unexplored. The theoretical connection between SSMs and classical adaptive filters makes this a natural research direction, but empirical validation is needed.

    \item \textbf{Gap 5: AI relay-MIMO equalization interaction.} While MIMO equalization (ZF, MMSE, SIC) and AI-based processing have been studied independently, their \textbf{joint} effect in a relay pipeline : where the relay denoises the signal before equalization separates the spatial streams : has not been investigated. Whether the gains from better relay processing and better equalization are additive, synergistic, or partially redundant is unknown.
\end{enumerate}

The following table summarizes the positioning of this work relative to the literature:

\begin{longtable}[]{@{}
  >{\raggedright\arraybackslash}p{(\linewidth - 4\tabcolsep) * \real{0.3333}}
  >{\raggedright\arraybackslash}p{(\linewidth - 4\tabcolsep) * \real{0.3333}}
  >{\raggedright\arraybackslash}p{(\linewidth - 4\tabcolsep) * \real{0.3333}}@{}}
\caption{Research positioning: comparison of this thesis against prior work in deep-learning-based relay communication.}\label{tbl:table17}\tabularnewline
\toprule\noalign{}
\begin{minipage}[b]{\linewidth}\raggedright
Aspect
\end{minipage} & \begin{minipage}[b]{\linewidth}\raggedright
Prior Work
\end{minipage} & \begin{minipage}[b]{\linewidth}\raggedright
This Thesis
\end{minipage} \\
\midrule\noalign{}
\endfirsthead
\toprule\noalign{}
\begin{minipage}[b]{\linewidth}\raggedright
Aspect
\end{minipage} & \begin{minipage}[b]{\linewidth}\raggedright
Prior Work
\end{minipage} & \begin{minipage}[b]{\linewidth}\raggedright
This Thesis
\end{minipage} \\
\midrule\noalign{}
\endhead
\bottomrule\noalign{}
\endlastfoot
Relay strategies compared & 1-2 (typically AF vs.~DF, or one AI method) & 9 (2 classical + 7 AI) \\
AI paradigms & Single (e.g., supervised only) & 4 (supervised, generative, adversarial, sequential) \\
Channel models & 1-2 (AWGN, sometimes Rayleigh) & 6 (AWGN + Rayleigh + Rician SISO; MIMO ZF/MMSE/SIC) \\
Complexity analysis & None & Normalized 3K comparison + inverted-U study \\
State space models & None for relay/physical-layer & Mamba $S_6$, Mamba-2 SSD, with crossover benchmark \\
Statistical rigor & Often missing & Wilcoxon test, 95\% CI, 10 trials per point \\
\end{longtable}



\section{Research Hypotheses}\label{sec:research-hypotheses}

Based on the theoretical analysis in Section \ref{sec:literature-review}, this thesis tests the following hypotheses:

\begin{enumerate}
    \item \textbf{H1 (Selective AI advantage at low SNR).} Some AI-based relay strategies achieve lower BER than AF and, on selected channels, lower BER than DF at low SNR (0--4 dB), where the non-linear denoising learned by the neural network may provide an advantage over AF's noise amplification and DF's hard-decision errors.

\emph{Rationale:} At low SNR, the Bayes-optimal relay function \(f^*(y) = \tanh(y/\sigma^2)\) is a smooth non-linear mapping that differs substantially from both AF's linear scaling and DF's hard sign function. A neural network can approximate \(f^*\) closely, while the classical methods cannot.

\item \textbf{H2 (DF dominance at high SNR).} The DF relay achieves the lowest BER at medium-to-high SNR (\(\geq 6\) dB), outperforming all AI methods.

\emph{Rationale:} At high SNR, \(f^*(y) \approx \text{sign}(y)\), which is exactly the DF operation. AI relays introduce a small but non-zero approximation error, so DF should be optimal in this regime.

\item \textbf{H3 (Inverted-U complexity curve).} There exists an optimal model size for relay denoising beyond which performance degrades due to overfitting. Specifically, models with $\sim$(100–200 parameters achieve performance comparable to models with 10–100$)\times$ more parameters.

\emph{Rationale:} The bias-variance analysis (Section \ref{sec:theoretical-basis-universal-approximation-and-denoising}) predicts that the low-complexity target function \(f^*\) requires minimal model capacity. Excess capacity increases variance without reducing bias.

\item \textbf{H4 (Architecture convergence at equal scale).} When all AI models are normalized to the same parameter count ($\sim$3,000), the performance differences between architectures narrow significantly, indicating that parameter count is a more important factor than architectural choice.

\emph{Rationale:} All architectures are universal approximators and the target function is simple. At equal capacity, architectural inductive biases provide diminishing returns.

\item \textbf{H5 (SSM speed advantage at long context).} Mamba-2 SSD trains significantly faster than Mamba S6 at longer context lengths (\(n \gg 11\)) due to chunk-parallel computation, while S6 is faster at the short relay window (\(n = 11\)).

\emph{Rationale:} The crossover between sequential S6 (\(O(n)\) serial steps) and chunk-parallel SSD (\(O(n/L)\) parallel matmuls of size \(L\)) depends on the ratio of kernel-launch overhead to arithmetic cost.

\item \textbf{H6 (Equalization gains are additive to relay gains).} The BER improvement from better equalization (ZF \(\to\) MMSE \(\to\) SIC) and the improvement from better relay processing (AF \(\to\) AI) are approximately additive in the dB domain.

\emph{Rationale:} The relay operates on Hop 1 (denoising) and the equalizer operates on Hop 2 (stream separation). Since these are independent processing stages, their contributions to BER reduction should be approximately independent.
\end{enumerate}

\section{Specific Objectives}\label{sec:specific-objectives}

\begin{enumerate}
\def\labelenumi{\arabic{enumi}.}
\item
  \textbf{Implement and compare nine relay strategies} spanning four learning paradigms: no learning (AF, DF), supervised learning (MLP, Hybrid), generative modeling (VAE, CGAN), and sequence modeling (Transformer, Mamba S6, Mamba-2 SSD).
\item
  \textbf{Evaluate across six channel/topology configurations:} AWGN, Rayleigh fading, and Rician fading (K=3) channels in SISO topology, and 2×2 MIMO Rayleigh with ZF, MMSE, and SIC equalization.
\item
  \textbf{Investigate the complexity--performance trade-off} by testing model architectures ranging from 0 parameters (classical) to 26,179 parameters (Mamba-2 SSD), and by conducting a normalized comparison at approximately 3,000 parameters.
\item
  \textbf{Determine whether state space models outperform attention mechanisms} for relay signal processing by comparing Mamba S6 and Transformer architectures at both original and normalized parameter counts, and by benchmarking at extended context lengths.
\item
  \textbf{Identify practical deployment recommendations} for selecting the appropriate relay strategy given specific operational constraints (SNR range, computational budget, channel environment).
\end{enumerate}

\section{Scope and limitations}\label{sec:scope-and-delimitations}

The following limitations define the scope of this study:

\begin{itemize}
\tightlist
\item
  \textbf{Modulation:} Primary experiments use BPSK. Extension experiments with QPSK and 16-QAM are presented in Section \ref{sec:higher-order-modulation-scalability-constellation-aware-training}, and 16-PSK in Section \ref{sec:input-normalization-and-csi-injection}, to evaluate hypothesis generalisability; however, 64-QAM and higher-order constellations are deferred to future work.
\item
  \textbf{Channel knowledge:} Perfect CSI is assumed at the receiver. Channel estimation errors are not modeled.
\item
  \textbf{Relay topology:} Single relay, two-hop, half-duplex. Multi-relay and full-duplex configurations are excluded.
\item
  \textbf{MIMO configuration:} \(2 \times 2\) spatial multiplexing with Rayleigh fading. Larger arrays and beamforming are not considered.
\item
  \textbf{Training regime:} Offline training on synthetic data. Online adaptation is not implemented.
\end{itemize}

These delimitations are chosen to enable a clean comparison of relay processing functions in isolation, without confounding factors from protocol-level or system-level complexity.

